\documentclass[11pt]{article}

\usepackage[a4paper,margin=1in]{geometry}
\usepackage{amsmath}
\usepackage{amssymb}
\usepackage{amsthm}
\usepackage{mathtools}
\usepackage[hidelinks]{hyperref}
\usepackage{doi}

\newtheorem{theorem}{Theorem}
\newtheorem{proposition}{Proposition}
\newtheorem{lemma}{Lemma}
\newtheorem{corollary}{Corollary}
\theoremstyle{definition}
\newtheorem{definition}{Definition}
\newtheorem{remark}{Remark}

\newcommand{\Heis}{\mathrm{Heis}}
\newcommand{\HS}{\mathrm{HS}}
\newcommand{\tr}{\operatorname{tr}}
\newcommand{\Cgen}{C^3_{\mathrm{gen}}}
\newcommand{\EPi}{E_\Pi}
\newcommand{\PiS}{\Pi_S}
\newcommand{\Deff}{\mathcal D}

\title{The Projective Residue as a Schur Complement:\\
Reduction of the Generation Split to the Chiral Asymmetry\\
of Projection Locking}
\author{J. Beau\\ Independent Researcher, France\\ \texttt{jerome.beau@cosmochrony.org}}
\date{2026}

\begin{document}
  \maketitle

  \begin{abstract}
    The fermionic sub-programme of Cosmochrony locates the three-generation mass split in the
    $J_3$-odd part of the squared projective endomorphism $E_\Pi^2$ restricted to the gauge-singlet
    generation triplet $\Cgen$, parametrised by a single real number $u$ through
    $E_\Pi^2|_{\Cgen}=\mathrm{diag}(1,\tfrac12+u,\tfrac12-u)$, with the even sector
    $\mathrm{diag}(1,\tfrac12,\tfrac12)$ already closed by Born--Infeld parity~\cite{Beau2026a34}.
    This note fixes the structural status of $u$ before any explicit construction of $E_\Pi$.
    First, the projected Dirac square admits a universal Feshbach/Schur form
    $E_\Pi=-\PiS\,D\,(1-P)\,D\,\PiS^{*}=-M^{\dagger}M$ with $P=\PiS^{*}\PiS$, exhibiting $E_\Pi$ as the
    Schur complement of the spinorial directions eliminated by the non-injective projection, negative
    semi-definite and vanishing in the injective limit.
    Second, the chiral block decomposition shows that $u$ is controlled by the $\Deff^{\pm}$-transported,
    generation-projected part of the antiunitary chiral equivariance defect
    \[
      \Delta_\chi(P)=\pi_{LL}-\tau\,\overline{\pi_{RR}}\,\tau^{-1}
    \]
    of the eliminated block $1-P$, rather than by a naive block difference.
    Third, the minimal non-injectivity $c\leftrightarrow q-c$ is chirally symmetric, so $u=0$ at the
    level of axioms A1--A3; a non-zero $u$ requires a chiral symmetry-breaking that can originate only
    in the projection-locking axiom A4.
    The Seeley--DeWitt conventions are locked so that the operator-level and spectral-level definitions
    of $u$ coincide on the flat effective metric, and a no-go result establishes that $u$ is not
    accessible to any finite front observable because chirality is a Lorentzian rather than a finite-fibre
    datum.
    The remaining open deliverable is reduced to three sharp questions about A4.
  \end{abstract}

  \noindent\textbf{Keywords:} projective endomorphism; Schur complement; non-injective projection;
  chiral asymmetry; projection locking; Born--Infeld saturation; generation split; Seeley--DeWitt
  expansion; admissible spinor bundle.

  \clearpage
  \tableofcontents

  \clearpage
  \section{Position and purpose}
    \label{sec:position}

    The fermionic sub-programme~\cite{Beau2026q14} derives the admissible spinor bundle $\mathcal S_\Pi$,
    the $V-A$ chiral structure, hypercharge rigidity, and a gauge-singlet three-generation factor
    $\Cgen$ as the spinorial face of the admissible Weil module after the Lorentzian complexification of
    its metaplectic lift.
    The generation mass split is carried by the projective endomorphism
    $E_\Pi\in\Gamma(\mathrm{End}(\mathcal S_\Pi))$, the zero-order spectral residue of non-injective
    projection in the spinorial sector~\cite{Beau2026q14}.
    On the rank-three admissible triplet $\Cgen=\mathrm{span}(e_0,e_+,e_-)$ with Cartan generator
    $J_3=\mathrm{diag}(0,1,-1)$, the squared residue takes the normal form
    \begin{equation}
      \label{eq:normalform}
      E_\Pi^2\big|_{\Cgen}=\frac{C_2-J_3^2}{C_2}+u\,J_3+v\,R_{\mathrm{mix}},
    \end{equation}
    where $C_2=2$ is the spin-one Casimir, $R_{\mathrm{mix}}$ is the off-diagonal generation mixing, and
    $u$ is the single real coefficient governing the split.
    The CP-even/real sector imposes $v=0$, and the even part is closed analytically by Born--Infeld
    parity on the equatorial admissible fibre~\cite{Beau2026a34}:
    \begin{equation}
      \label{eq:even}
      (E_\Pi^2)_{\mathrm{even}}\big|_{\Cgen}=\frac{C_2-J_3^2}{C_2}=\mathrm{diag}\!\left(1,\tfrac12,\tfrac12\right).
    \end{equation}
    The only remaining unknown is therefore
    \begin{equation}
      \label{eq:udef}
      u=\frac{\big\langle \Pi_{J_\Pi\text{-odd}}(E_\Pi^2),\,J_3\big\rangle_{\HS}}{\langle J_3,J_3\rangle_{\HS}}
      =\tfrac12\big(\langle e_+|E_\Pi^2|e_+\rangle-\langle e_-|E_\Pi^2|e_-\rangle\big).
    \end{equation}
    This note records what is now established about $u$ and isolates the single conceptual obstruction
    that remains, so that the explicit construction of $E_\Pi$ can proceed in a locked frame.
    The companion front observables of the angular and oriented-frontier
    notes~\cite{Beau2026a34} are not the locus of $u$; they explain the obstructions and provide null
    controls, but $u$ lives in the zero-order residue of the projected Dirac operator.
    The numerical front controls are recorded in Appendix~\ref{sec:controls}.

  \section{The projective residue as a Schur complement}
    \label{sec:schur}

    Let $\PiS:\mathcal S\to\mathcal S_\Pi$ be the restriction of the admissible projection to the
    spinorial bundle, $\PiS^{*}$ the admissible minimal lift, and
    \begin{equation}
      \label{eq:dpidef}
      D_{\Pi,g,A}:=\PiS\circ D_{g,A}\circ \PiS^{*},\qquad P:=\PiS^{*}\PiS,
    \end{equation}
    following~\cite{Beau2026q14}.
    One does not assume $\PiS\PiS^{*}=\mathrm{id}_{\mathcal S_\Pi}$; this holds only in the injectively
    resolved limit.
    Symbol-compatibility means the projection defect affects only the zero-order part of
    $D^2_{\Pi,g,A}$, leaving the Clifford principal symbol intact.

    \begin{proposition}[Schur form of the projective residue]
      \label{prop:schur}
      Under symbol-compatibility, the projective endomorphism of~\cite{Beau2026q14} satisfies
      \begin{equation}
        \label{eq:schur}
        \EPi=-\PiS\,D_{g,A}\,(1-P)\,D_{g,A}\,\PiS^{*}=-M^{\dagger}M,
        \qquad M:=(1-P)\,D_{g,A}\,\PiS^{*}.
      \end{equation}
      In particular $\EPi\preceq 0$, it is zero-order, and $\EPi=0$ if and only if $1-P=0$, i.e.\ in the
      injective limit.
    \end{proposition}

    \begin{proof}
      With $P=\PiS^{*}\PiS$ idempotent, $D^2_{\Pi,g,A}=\PiS D_{g,A}P D_{g,A}\PiS^{*}$.
      The uncompressed square $\PiS D_{g,A}^2\PiS^{*}$ reconstructs the effective spin--gauge Lichnerowicz
      operator on $\mathcal S_\Pi$ up to symbol, so that
      \begin{equation*}
        \mathrm{Lich}(g,A)\approx \PiS D_{g,A}^2\PiS^{*}=D^2_{\Pi,g,A}+\PiS D_{g,A}(1-P)D_{g,A}\PiS^{*}.
      \end{equation*}
      Subtracting the Lichnerowicz block from $D^2_{\Pi,g,A}$ gives~\eqref{eq:schur}.
      Since $1-P$ is a self-adjoint projector and $D_{g,A}$ is self-adjoint,
      $\PiS D_{g,A}(1-P)D_{g,A}\PiS^{*}=M^{\dagger}M\succeq 0$, hence $\EPi=-M^{\dagger}M\preceq 0$.
      The first-order part vanishes by symbol-compatibility, leaving a zero-order endomorphism, and
      $\EPi=0\Leftrightarrow M=0\Leftrightarrow(1-P)D_{g,A}\PiS^{*}=0$, which for a faithful symbol holds
      exactly when $1-P=0$.
    \end{proof}

    Proposition~\ref{prop:schur} makes precise the statement of~\cite{Beau2026q14} that $E_\Pi$ is not a
    postulated finite Dirac datum but is induced by the non-injective projection itself: it is the Schur
    complement of the eliminated sector $1-P$.

  \section{Chiral block reduction and the localisation of $u$}
    \label{sec:chiral}

    The chirality operator $\gamma_5=i\gamma^0\gamma^1\gamma^2\gamma^3$ of the effective Clifford algebra
    exists as a global section because of the Lorentzian closure of the effective
    metric~\cite{Beau2026q11}.
    The Dirac operator is chirality-odd, hence off-diagonal in the chiral splitting
    $\mathcal S_\Pi=\mathcal S_L\oplus\mathcal S_R$:
    \begin{equation}
      \label{eq:blocks}
      D_{g,A}=\begin{pmatrix}0&\Deff^{-}\\ \Deff^{+}&0\end{pmatrix},
      \qquad
      1-P=\begin{pmatrix}\pi_{LL}&\pi_{LR}\\ \pi_{RL}&\pi_{RR}\end{pmatrix},
    \end{equation}
    with $\pi_{LL},\pi_{RR}$ self-adjoint and $\pi_{RL}=\pi_{LR}^{\dagger}$.

    \begin{proposition}[Chiral diagonal blocks of the residue]
      \label{prop:chiralblocks}
      With the conventions of~\eqref{eq:blocks},
      \begin{equation}
        \label{eq:diag}
        E_{LL}=-\PiS\,\Deff^{-}\pi_{RR}\Deff^{+}\,\PiS^{*},\qquad
        E_{RR}=-\PiS\,\Deff^{+}\pi_{LL}\Deff^{-}\,\PiS^{*}.
      \end{equation}
      The left-admissibility condition $P_R\EPi P_R=0$ of~\cite{Beau2026q14} is equivalent to $\pi_{LL}=0$:
      the non-injective projection eliminates no purely left-handed direction.
    \end{proposition}

    \begin{proof}
      Direct multiplication of~\eqref{eq:blocks} gives
      \begin{equation*}
        D_{g,A}(1-P)D_{g,A}=
        \begin{pmatrix}\Deff^{-}\pi_{RR}\Deff^{+}&\Deff^{-}\pi_{RL}\Deff^{-}\\
          \Deff^{+}\pi_{LR}\Deff^{+}&\Deff^{+}\pi_{LL}\Deff^{-}\end{pmatrix}.
      \end{equation*}
      Conjugating by $\PiS,\PiS^{*}$ and using $\EPi=-\PiS D_{g,A}(1-P)D_{g,A}\PiS^{*}$
      gives~\eqref{eq:diag}.
      The right-right block of $\EPi$ is $E_{RR}=-\PiS\Deff^{+}\pi_{LL}\Deff^{-}\PiS^{*}$, which vanishes
      for a faithful Weyl symbol exactly when $\pi_{LL}=0$.
    \end{proof}

    \begin{corollary}[Localisation of $u$]
      \label{cor:localise}
      After restriction to $\Cgen$, the $J_3$-odd part of $E_\Pi^2$ is controlled by the chiral asymmetry
      of the eliminated block,
      \begin{equation}
        \label{eq:uasym}
        u\ \propto\ \pi_{RR}-\pi_{LL},
      \end{equation}
      dressed by the Weyl transport $\Deff^{\pm}$.
      The split $u$ is therefore an intrinsic spectral asymmetry of the chiral part of the projective
      residue, not a shell average of any front cocycle.
    \end{corollary}

  \section{Stratification: $u$ vanishes unless A4 breaks chiral symmetry}
    \label{sec:strat}

    The minimal non-injectivity is the parity involution $c\leftrightarrow q-c$ on conjugate Weil
    blocks~\cite{Beau2026fe}.
    Its spinorial lift is the antilinear Born--Infeld parity $J_\Pi$, which exchanges the chiral
    subbundles, $J_\Pi P_L=P_R J_\Pi$~\cite{Beau2026q14}.
    A parity that exchanges $\mathcal S_L\leftrightarrow\mathcal S_R$ treats the two chiralities
    symmetrically, so the minimal eliminated block is chirally symmetric,
    $\pi_{LL}=\pi_{RR}$.

    \begin{proposition}[Chiral-symmetry stratification of $u$]
      \label{prop:strat}
      If $\pi_{LL}=\pi_{RR}$ then $u=0$.
      The conjugate pair $\{V^{(c)},V^{(q-c)}\}$ carries no preferred orientation at the level of axioms
      A1--A3, consistently with the non-premature-selection axiom A3~\cite{Beau2026m}.
      Hence at the A1--A3 level $u=0$.
      A non-zero $u$ requires $\pi_{LL}\neq\pi_{RR}$, a chiral symmetry-breaking that can originate only in
      the projection-locking axiom A4.
    \end{proposition}

    \begin{proof}
      By Corollary~\ref{cor:localise}, $u\propto\pi_{RR}-\pi_{LL}$, so $\pi_{LL}=\pi_{RR}$ forces $u=0$.
      The minimal non-injectivity is the chirally symmetric parity $c\leftrightarrow q-c$, and A3 forbids
      selection among open directions before locking, so no chiral orientation is fixed by A1--A3.
      The only axiom that selects a resolved direction is A4, ``resolution occurs only when continuous
      Born--Infeld saturation meets the discrete shell support''~\cite{Beau2026m}; any chiral asymmetry
      $\pi_{LL}\neq\pi_{RR}$ must therefore arise there.
    \end{proof}

    The left-admissibility branch $\pi_{LL}=0$~\cite{Beau2026q14} is thus the maximal chiral asymmetry
    compatible with A4 locking, and the sign of $u$ is the spontaneous $V-A$ choice exchanged under
    $P_L\leftrightarrow P_R$.
    The even sector~\eqref{eq:even} is the chirally symmetric part of the same Born--Infeld
    saturation~\cite{Beau2026a34, Beau2026c}; $u$ is the chirally asymmetric part of that same locking.
    This is the structural sense in which the Higgs amplitude mode of the saturated projective condensate
    and the generation split are read from a single A4 event.

  \section{Spinorial locking and the $J_\Pi$-equivariance defect}
    \label{sec:spinorial-locking}

    This section refines Proposition~\ref{prop:strat}: the chirally symmetric branch is the transported
    equivalence $\pi_{LL}=\tau\,\overline{\pi_{RR}}\,\tau^{-1}$, not the naive block equality
    $\pi_{LL}=\pi_{RR}$, the bar being the complex conjugation carried by the antiunitary Born--Infeld
    lift $J_\Pi$.

    Let $V=\mathrm{span}(v_+,v_-)$ denote the fundamental spinorial doublet, with Cartan generator
    $Hv_\pm=\pm v_\pm$.
    The spinorial Born--Infeld lift of~\cite{Beau2026q14} is the antiunitary structure
    \begin{equation}
      \label{eq:jpidef}
      J_\Pi=\epsilon\circ\mathrm{conj},\qquad \epsilon=\sigma_z\sigma_x=i\sigma_y,
    \end{equation}
    so that, in the standard convention, $J_\Pi v_+=-v_-$ and $J_\Pi v_-=v_+$.
    Thus $J_\Pi^2=-1$ on $V$: this is the pseudoreal, or quaternionic, structure of the
    fundamental spinorial sector.

    On the symmetric square $\mathrm{Sym}^2(V)=\mathrm{span}(e_0,e_+,e_-)$ with $e_0=\sqrt{2}\,v_+v_-$,
    $e_+=v_+^2$, $e_-=v_-^2$, the induced antiunitary $J_\Pi^{(2)}$ acts by
    \begin{equation}
      \label{eq:jpisym}
      J_\Pi^{(2)}e_0=-e_0,\qquad J_\Pi^{(2)}e_+=e_-,\qquad J_\Pi^{(2)}e_-=e_+.
    \end{equation}
    Consequently $(J_\Pi^{(2)})^2=+1$ on $\mathrm{Sym}^2(V)$: the generation triplet carries a real
    involution, although the underlying doublet is pseudoreal.
    With $J_3=\mathrm{diag}(0,1,-1)$ on $(e_0,e_+,e_-)$, the formulas~\eqref{eq:jpisym} give
    \begin{equation}
      \label{eq:jpij3}
      J_\Pi^{(2)}J_3(J_\Pi^{(2)})^{-1}=-J_3,
    \end{equation}
    so the Born--Infeld spinorial conjugation exchanges the two outer generation weights.

    We now apply this to the Schur residue.
    Let $\tau:\mathcal S_R\to\mathcal S_L$ be the unitary intertwiner realising the chirality exchange
    $J_\Pi P_L=P_R J_\Pi$ of~\cite{Beau2026q14} in the chiral frame.
    Equivariance of the locking projector, $[P,J_\Pi]=0$, then forces on the diagonal chiral blocks of
    $1-P$ not the naive equality but the transported, complex-conjugated equivalence
    \begin{equation}
      \label{eq:transported}
      \pi_{LL}=\tau\,\overline{\pi_{RR}}\,\tau^{-1}.
    \end{equation}
    The relevant obstruction is consequently the antiunitary defect
    \begin{equation}
      \label{eq:defect}
      \Delta_\chi(P):=\pi_{LL}-\tau\,\overline{\pi_{RR}}\,\tau^{-1},
    \end{equation}
    which reduces to $\pi_{LL}-\tau\pi_{RR}\tau^{-1}$ on the real structure fixed by $J_\Pi^{(2)}$ on the
    generation triplet.

    \begin{lemma}[Chiral obstruction of spinorial locking]
      \label{lem:chiral-obstruction}
      Assume the projective Dirac transport is $J_\Pi$-compatible, $J_\Pi D_{g,A}J_\Pi^{-1}=\pm D_{g,A}$ and
      $J_\Pi\PiS J_\Pi^{-1}=\PiS$, as in the spinorial Born--Infeld lift of~\cite{Beau2026q14}.
      If the spinorial locking projector is $J_\Pi$-equivariant, $[P,J_\Pi]=0$, then the Schur residue
      $\EPi=-\PiS D_{g,A}(1-P)D_{g,A}\PiS^{*}$ satisfies $J_\Pi\EPi^2 J_\Pi^{-1}=\EPi^2$.
      Consequently $\langle e_+|\EPi^2|e_+\rangle=\langle e_-|\EPi^2|e_-\rangle$, and the split coefficient
      $u$ of~\eqref{eq:udef} vanishes.
      Hence a non-zero generation split requires a failure of $J_\Pi$-equivariance at projection locking,
      \begin{equation*}
        u\neq0\quad\Longrightarrow\quad [P,J_\Pi]\neq0.
      \end{equation*}
    \end{lemma}

    \begin{proof}
      Since $J_\Pi^{(2)}e_+=e_-$ and~\eqref{eq:jpij3} holds, the two outer generation weights are conjugate
      under the Born--Infeld spinorial involution.
      The assumptions $J_\Pi D_{g,A}J_\Pi^{-1}=\pm D_{g,A}$, $J_\Pi\PiS J_\Pi^{-1}=\PiS$ and $[P,J_\Pi]=0$
      imply $J_\Pi\EPi J_\Pi^{-1}=\pm\EPi$, up to the irrelevant sign induced by the chirality-odd Dirac
      transport.
      Squaring removes this sign and gives $J_\Pi\EPi^2 J_\Pi^{-1}=\EPi^2$.
      Because $J_\Pi$ is antiunitary,
      $\langle J_\Pi\psi|\EPi^2|J_\Pi\psi\rangle=\overline{\langle\psi|\EPi^2|\psi\rangle}$, and $\EPi^2$ is
      self-adjoint, so this expectation value is real.
      Taking $\psi=e_+$ and using $J_\Pi e_+=e_-$ gives
      $\langle e_-|\EPi^2|e_-\rangle=\langle e_+|\EPi^2|e_+\rangle$.
      The definition~\eqref{eq:udef} of $u$ then gives $u=0$, and the final implication is the
      contrapositive.
    \end{proof}

    \begin{remark}[Necessary condition, not a sufficient one]
      \label{rem:necessary}
      The condition $[P,J_\Pi]\neq0$ is necessary for $u\neq0$, but not sufficient.
      The commutator may be supported outside the generation triplet, or in components that do not survive
      the Dirac transport and the final projection to $\Cgen$.
      The magnitude of $u$ is therefore not determined by the abstract commutator itself: it is set by the
      $\Deff^{\pm}$-transported, generation-projected part of the antiunitary defect $\Delta_\chi(P)$
      of~\eqref{eq:defect}.
      Determining this transported and projected component is the remaining A4-level deliverable.
    \end{remark}

  \section{Locked Seeley--DeWitt conventions: operator $u$ equals spectral $u$}
    \label{sec:conventions}

    The projected square is Laplace-type, $D^2_{\Pi,g,A}=-(\nabla^{S,A})^2+E$, with
    $E=\tfrac{R}{4}+\tfrac12\gamma^\mu\gamma^\nu F^a_{\mu\nu}T^a+\EPi$~\cite{Beau2026q14}, and the
    $\gamma_5$-weighted Seeley--DeWitt coefficient $a_4=(4\pi)^{-2}\tr(\tfrac12E^2-\tfrac16 RE+\dots)$ in
    the convention of~\cite{Beau2026q14} consistent with the heat-kernel manual.
    The bosonic synthesis fixes the spectral normalisation: $a_2$ is the horizontal gravitational
    response and $a_4$ the vertical gauge response of the spectral entropy functional, with induced
    couplings carrying a common factor at each Seeley--DeWitt order~\cite{Beau2026q13, Beau2026q12}.

    Three facts make the operator-level $u$ of~\eqref{eq:udef} coincide with the spectral coefficient
    that enters the effective amplitude.
    First, $u$ is quadratic in $\EPi$, so it is immune to the linear sign convention of $E$.
    Second, on the gauge-singlet $\Cgen$ the gauge curvature term vanishes, and on the flat effective
    metric $g^{\mu\nu}=2\eta^{\mu\nu}$~\cite{Beau2026q11} all curvature cross-terms vanish, in particular
    $-\tfrac16 R\EPi=0$; only $\tfrac12\EPi^2$ survives, so no linear-in-$\EPi$ term can separate the
    operator and spectral readings.
    Third, the common $a_4$ normalisation factor cancels in the generation ratio
    $(\tfrac12+u)/(\tfrac12-u)$, so $u$ is independent of the spectral normalisation while the absolute
    mass scale carries it.
    A curvature correction $-\tfrac16 R\EPi$ would reappear only on a curved background, as a
    gravitational correction to masses, and does not affect the structural $u$.

  \section{Chirality is not a finite-fibre label}
    \label{sec:nogo}

    The chirality $\gamma_5$ is the chirality operator of the effective spacetime Clifford
    algebra~\cite{Beau2026q14}, existing because of the Lorentzian closure
    $g^{\mu\nu}=2\eta^{\mu\nu}$~\cite{Beau2026q11}.
    The finite Weil module selects only the compact pseudoreal form $\mathfrak{su}(2)$, in which the
    fundamental is self-conjugate~\cite{Beau2026a27}; the Lorentzian complexification is what removes the
    equivalence and renders $\mathcal S_L$ and $\mathcal S_R$ inequivalent~\cite{Beau2026q14}.

    \begin{remark}[No finite chiral edge label]
      \label{rem:nogo}
      A cascade edge $g\to gs$ is a Heisenberg translation, not a symplectic element, so it carries no
      canonical metaplectic lift; and the chiral eigendirections are the complex combinations of the real
      generators, so a real edge has no definite chirality.
      On the admissible front the only orientation-odd coordinate datum reduces to the sign of the
      cocycle, which is the eliminated-amplitude proxy, not a chirality.
      Consequently $u$ cannot be extracted from any finite front observable: the chiral doubling
      $\mathcal S_L\oplus\mathcal S_R$ and the zero-order Dirac residue are required.
      This explains, structurally, every vanishing of the symmetric front and shell observables: each
      lacks the chiral asymmetry that lives in the Lorentzian residue.
    \end{remark}

  \section{Status of results}
    \label{sec:status}

    \noindent\emph{Proved / structural (this note).}
    \begin{itemize}
      \item $\EPi=-\PiS D(1-P)D\PiS^{*}=-M^{\dagger}M$ is the Schur complement of the eliminated sector,
      negative semi-definite, zero-order, vanishing in the injective limit (Proposition~\ref{prop:schur}).
      \item Chiral diagonal blocks~\eqref{eq:diag}; left-admissibility $\Leftrightarrow\pi_{LL}=0$
      (Proposition~\ref{prop:chiralblocks}).
      \item $u\propto\pi_{RR}-\pi_{LL}$ on $\Cgen$ (Corollary~\ref{cor:localise}).
      \item $\pi_{LL}=\pi_{RR}\Rightarrow u=0$; non-zero $u$ requires A4 chiral symmetry-breaking
      (Proposition~\ref{prop:strat}).
      \item The symmetric-square action of the Born--Infeld involution is fixed: $J_\Pi^2=-1$ on the
      fundamental doublet, $(J_\Pi^{(2)})^2=+1$ on the generation triplet, and it anticommutes with
      $J_3$~\eqref{eq:jpij3}.
      Hence $u\neq0$ requires a failure of $J_\Pi$-equivariance at projection locking
      (Lemma~\ref{lem:chiral-obstruction}); the magnitude of $u$ depends on the $\Deff^{\pm}$-transported,
      generation-projected component of the antiunitary defect
      $\Delta_\chi(P)=\pi_{LL}-\tau\,\overline{\pi_{RR}}\,\tau^{-1}$~\eqref{eq:defect}.
      \item Operator $u$ equals spectral $u$ on the flat effective metric; $u$ normalisation-independent in
      the ratio (Section~\ref{sec:conventions}).
      \item Chirality is Lorentzian, not a finite-fibre label; $u$ is not a front observable
      (Remark~\ref{rem:nogo}).
    \end{itemize}

    \noindent\emph{Conditional.}
    The left-admissibility branch $\pi_{LL}=0$~\cite{Beau2026q14} is an admissible branch selected by the
    orientation-compatible assumption; its derivation from A4 is open.

    \noindent\emph{Open.}
    The magnitude $|u|$ as the spectral weight of the chiral block eliminated by A4 locking.

  \section{Open deliverables}
    \label{sec:open}

    The conceptual obstruction is reduced to three questions about the projection-locking axiom A4 and
    the Born--Infeld saturation that defines it~\cite{Beau2026c, Beau2026m}.
    \begin{enumerate}
      \item Does A4 impose a chiral asymmetry $\pi_{RR}\neq\pi_{LL}$, or only an unoriented discrete
      selection? If only the latter, left-admissibility remains a conditional branch.
      \item If A4 imposes an asymmetry, does it fix only its sign (the spontaneous $V-A$ choice), or also
      its magnitude?
      \item Is the magnitude tied to the degree of Born--Infeld saturation, to the discrete shell support,
      or to a still-open normalisation?
    \end{enumerate}
    These are the last conceptual locks before the explicit construction of $E_\Pi$ as the Schur
    complement~\eqref{eq:schur} with an explicit eliminated block $1-P$.

  \section{Conclusion}
    \label{sec:conclusion}

    The generation split $u$ is now localised with no remaining ambiguity of method: it is the chiral
    asymmetry $\pi_{RR}-\pi_{LL}$ of the spinorial directions eliminated by the non-injective projection,
    read off the Schur complement $\EPi=-M^{\dagger}M$, projected onto the gauge-singlet triplet $\Cgen$.
    It vanishes identically for the chirally symmetric minimal non-injectivity, and is non-zero only
    through a chiral symmetry-breaking carried by the projection-locking axiom A4.
    The polarisation, the Seeley--DeWitt conventions, and the operator/spectral identification are locked;
    the front observables are confirmed to be null controls rather than carriers of $u$.
    The single remaining question is whether A4 derives the chiral asymmetry of the eliminated block, and
    with what magnitude, which is the next text to settle from the Born--Infeld saturation
    structure~\cite{Beau2026c}.

    \appendix
  \section{Numerical front controls}
    \label{sec:controls}

    This appendix records the scalar front diagnostics obtained in the session runs of the
    oriented-frontier pipeline.
    They are the null controls referenced in Section~\ref{sec:nogo}: the scalar front explains the
    obstruction and bounds the maximal locking, but does not carry $u$, in accordance with
    Remark~\ref{rem:nogo}.
    The values below are from the session runs and are to be reconciled with the archived script output
    before any deposition.

    The oriented-frontier increment is the cocycle $\Delta A_c(g,s)=c\,(z(s)+a(g)s_b)$, reduced to
    $c\,a(g)s_b$ on admissible generators.
    Its symmetric front average vanishes identically,
    \begin{equation}
      \label{eq:thetascalar}
      \Theta_{\mathrm{scalar}}=\max_{e\in\partial^+ S_m}\big|\theta^+_{\mathrm{raw}}(e)\big|=0
      \qquad (q\in\{61,101,151\}),
    \end{equation}
    the signature of the $\varphi$-re-pairing of the oriented frontier.
    This is the null control, not a failure: it is the scalar shadow of the chiral asymmetry that lives
    in the Lorentzian residue (Section~\ref{sec:nogo}).

    The maximal locking bound, normalised by the angle quantum $2\pi/q$, is $q$-invariant.
    The bare bound $\Theta_{\max}=\langle|\Delta A_c|\rangle$ decays as $\sim 2.118/q$ (the $2\pi/q$
    convention), so the $q$-invariant content is the product:
    \begin{center}
      \begin{tabular}{cc}
        \hline
        $q$ & $\Theta_{\max}\cdot q$\\
        \hline
        $61$ & $2.1141$\\
        $101$ & $2.1221$\\
        $151$ & $2.1183$\\
        \hline
      \end{tabular}
    \end{center}
    giving $\Theta_{\max}\cdot q\approx 2.118$, stable to within $\pm 0.2\%$ across the tested primes.
    This is the $q$-stable maximal-locking bound of the scalar front; it is a diagnostic of the bounded
    transport, not the spectral split $u$.

    The pair-saturation and the BFS regime are
    \begin{equation}
      \label{eq:nsat}
      n_3^{\mathrm{obs}}=2,\qquad |\mathrm{nodes}|=21569 \ \text{at depth }14\ (\text{identical across }q),
    \end{equation}
    confirming that the measurement lives at shell depth $\le 2$ on the $q$-independent nilpotent ball,
    in the pre-wraparound regime ($q^3\gg 21569$).
    The saturation rank $n_3^{\mathrm{obs}}=2$ is the $\sigma_{\mathrm{pair}}$ saturation of the front and
    is distinct from the rank-three generation count $\sigma_c(n_3)=3$ of~\cite{Beau2026q14}, which is a
    separate admissible-rank observable.
    With only two contributing shells, the $q$-stability of~\eqref{eq:thetascalar}--\eqref{eq:nsat} is the
    operative robustness check; the constancy of any per-shell profile is a weak test on two points.

    \bibliographystyle{plain}
    \bibliography{cosmochrony-bibliography}

\end{document}
